\documentclass[11pt,a4paper]{article}

\usepackage[T1]{fontenc}
\usepackage[utf8]{inputenc}
\usepackage{lmodern}
\usepackage[margin=2.2cm,headheight=14pt]{geometry}
\usepackage{microtype}
\usepackage{xcolor}
\usepackage{booktabs}
\usepackage{tabularx}
\usepackage{longtable}
\usepackage{array}
\usepackage{enumitem}
\usepackage{caption}
\usepackage{graphicx}
\usepackage{fancyhdr}
\usepackage[hidelinks]{hyperref}

\definecolor{accent}{HTML}{1C5CAB}
\definecolor{rule}{HTML}{C3C2B7}
\definecolor{soft}{HTML}{F3F6FA}
\definecolor{ink}{HTML}{16202C}

\setlength{\emergencystretch}{3em}
\setlist{itemsep=3pt,topsep=4pt,parsep=0pt}
\captionsetup{font=small,labelfont={bf,color=accent}}
\sloppy

\pagestyle{fancy}
\fancyhf{}
\fancyhead[L]{\small\color{ink}EC8202 Applied Big Data Engineering}
\fancyhead[R]{\small\color{ink}Fleet Operations Data System}
\fancyfoot[C]{\small\thepage}
\renewcommand{\headrulewidth}{0.4pt}

\title{\vspace{-1.1cm}\textbf{A Lambda Data Pipeline for\\Ride-Hailing Fleet Operations}}
\author{EC8202 Applied Big Data Engineering\\[8pt]
\normalsize
\begin{tabular}{cc}
\textbf{Jackshan Venujan G.S} & \textbf{Dinojan V.}\\
Register No. EG/2021/4566 & Register No. EG/2021/4487
\end{tabular}}
\date{}

\newcommand{\imageplaceholder}[3]{%
  \begin{figure}[htbp]\centering
  \IfFileExists{figures/#1}{%
    \includegraphics[width=#3\textwidth,height=0.58\textheight,keepaspectratio]{figures/#1}%
  }{%
    \makebox[\textwidth][c]{\fbox{\parbox[c][2.5cm][c]{0.90\textwidth}{%
      \centering\small\itshape Add the requested project image here.}}}%
  }%
  \caption{#2}\label{fig:#1}
  \end{figure}}

\begin{document}
\maketitle
\thispagestyle{fancy}

\vspace{-0.5cm}
\begin{center}
\begin{minipage}{0.94\textwidth}\small
\color{ink}\textbf{Summary.}
This project shows how a ride-hailing company can watch its fleet and review its daily costs.
It uses two data paths. The live path gives a quick view of vehicle activity. The daily path uses
the full day's vehicle data and a cost file to calculate each vehicle's income and profit or loss.
The project uses Kafka, Spark, Airflow, PostgreSQL, an API, and monitoring tools. It runs with
Docker Compose. One simulated day takes five real minutes.
\end{minipage}
\end{center}

\vspace{0.25cm}
\tableofcontents
\vspace{0.35cm}

\section{Use Case and Business Need}

A ride-hailing company needs to know how its vehicles are being used. Dispatchers need current
information to decide where to send vehicles. The finance team needs to know whether each vehicle
earned more than it cost to run.

These needs have different timing. The live view must update within seconds. It can be an estimate
because new events keep arriving. The daily result must use the full day of data and the cost file.
It is used to review vehicle performance, so the company needs a complete result.

\begin{table}[htbp]\centering\small
\caption{The two questions the system answers.}
\begin{tabularx}{\textwidth}{@{}p{3.0cm}XX@{}}
\toprule
\textbf{Need} & \textbf{Live fleet view} & \textbf{Daily review}\\
\midrule
Question & What are the vehicles doing now? & Which vehicles made or lost money today?\\
Data & Vehicle location, speed, status, and fares & Full vehicle history plus fuel and repair costs\\
When needed & Every few seconds & Once each simulated day\\
Result & Recent activity and alerts & Daily income, costs, and profit or loss\\
Accuracy & A quick estimate is useful & Use the complete available data\\
\bottomrule
\end{tabularx}
\end{table}

The system uses simulated vehicles and simulated costs. This lets the team show the full process
without connecting to real drivers, cars, fuel providers, or garages. The two simulated sources
describe the same set of vehicles.

\section{Architecture Choice}

The project uses a Lambda design. It has a live path and a daily path. Both paths use the same
vehicle events, but they answer different questions.

\subsection{Why two paths are needed}

The live path reads new vehicle events and updates the current fleet view. Spark processes the
events in small groups. This keeps the view fresh. Counts and earnings can be approximate because
some events may arrive late and the day is not yet complete.

The daily path waits for a full day's vehicle history and the cost file. It then recalculates the
day and joins vehicle income with fuel and repair costs. This gives the project its main daily
financial result.

\begin{table}[htbp]\centering\small
\caption{Why the project uses Lambda instead of one stream-only path.}
\begin{tabularx}{\textwidth}{@{}p{3.0cm}XX@{}}
\toprule
\textbf{Point} & \textbf{Lambda design used here} & \textbf{One stream-only path}\\
\midrule
Live answer & A fast view is produced from incoming events & Can also produce a fast view\\
Daily cost file & The file is joined during the daily job & The file would need special handling in the stream\\
Late events & The full history can be read again & Events past the allowed wait time may be missed\\
Daily result & The day is recalculated with its costs & A live estimate is not enough for the final daily review\\
\bottomrule
\end{tabularx}
\end{table}

A batch-only system would not help dispatchers during the day. A stream-only system would not give
the same full-day review after costs arrive. The two paths make the difference clear. The dashboard
labels the live view as an estimate and the daily view as the full-day result.

There is a cost to this design. The system keeps a live path and a daily path. Some parts of the
work are repeated. The daily path is still useful because the cost file arrives once a day and
needs to be checked against the full vehicle history.

\section{System Architecture}

\imageplaceholder{architecture.png}{End-to-end architecture. Blue is the speed path, orange is the batch path, and green is storage. The master dataset is written by the speed layer and read again by the batch layer. This is the defining loop of a Lambda architecture.}{0.95}

The vehicle simulator sends updates to Kafka. Spark Structured Streaming checks the events and
updates the live data in PostgreSQL. It also saves the full event history in Parquet. Invalid
events are kept for review, and the system creates alerts when a vehicle stays idle too long.

The cost simulator makes one file for each simulated day. Airflow waits until the file and the
matching vehicle history are ready. It then starts a Spark batch job and checks the result. The
batch job writes the daily result to PostgreSQL and creates report files. The API reads the data
and supplies it to the fleet dashboard.

Prometheus collects health measures from the long-running services. The short batch job sends its
measures through Pushgateway. Grafana displays the measures and alerts.

\section{Technology Choices}

The tools were chosen to match the two data sources and the work they must do.

\begin{table}[htbp]\centering\small
\caption{Tools and why they fit this fleet system.}
\begin{tabularx}{\textwidth}{@{}>{\raggedright\arraybackslash}p{2.8cm}>{\raggedright\arraybackslash}p{3.7cm}>{\raggedright\arraybackslash}X@{}}
\toprule
\textbf{Part} & \textbf{Tool} & \textbf{Reason for using it}\\
\midrule
Vehicle data & Kafka & It receives and holds the stream of vehicle updates. Spark can read from it while the simulator keeps sending events.\\
Live processing & Spark Structured Streaming & It checks incoming events, makes recent summaries, updates vehicle status, and raises idle alerts.\\
Daily processing & Spark batch & It reads the complete day of vehicle history and joins it with the daily cost file.\\
Daily coordination & Airflow & It waits for both inputs, starts the batch job, checks the result, and records that the day is complete.\\
History storage & Parquet & It stores a large history in a compact form. The batch job can read the data for one day.\\
Serving storage & PostgreSQL & It stores the live and daily results in tables that the API can query.\\
Dashboard API & FastAPI & It provides a simple way for the dashboard and other users to request results.\\
Monitoring & Prometheus and Grafana & Prometheus gathers health measures and checks alert rules. Grafana displays the pipeline status.\\
Local setup & Docker Compose & It starts the services together and keeps saved data in named volumes between ordinary stops.\\
\bottomrule
\end{tabularx}
\end{table}

Spark is used for both live and daily work. This keeps the main data processing in one tool. The
live job uses event time windows and can give a quick estimate. The daily job reads the full day
again and produces the daily result. Airflow coordinates the batch work, but Spark does the data
calculations.

\section{Data Sources and Simulated Time}

\subsection{Live vehicle updates}

The vehicle simulator represents twelve vehicles. Every two seconds it sends a vehicle update with
a location, speed, status, time, and fare. A vehicle can be idle, travelling to a passenger, or
carrying a passenger.

Some events are made invalid on purpose. Examples include a missing location, an impossible speed,
an unknown status, or a fare increase while the vehicle is parked. Spark checks each event. Invalid
events are kept apart so they can be reviewed instead of being used as good data.

\subsection{Daily cost information}

A second simulator creates one file for each simulated day. Each row may contain a vehicle's fuel
cost, maintenance cost, distance, and service note. Some vehicles are left out. Some reported
distances do not match the GPS distance. The daily job keeps these cases in the result and adds a
flag so they are visible.

\subsection{Clock used in the demo}

One simulated day takes five real minutes. The calendar uses Sri Lankan time. This speed lets the
team show a daily process in a short demo. A simulated hour takes about twelve and a half seconds.
The live job uses 10-second and 15-second processing intervals. The daily job runs about once every
five real minutes, after its inputs are ready.

The shared clock starts from one saved time so that all services use the same simulated date. This
prevents the vehicle source, cost source, and daily job from using different day boundaries.

\section{Processing and Daily Calculations}

\subsection{Live processing}

Kafka carries the vehicle messages to Spark Structured Streaming. Spark checks the message shape
and values, then adds a zone based on the vehicle location. It writes each event to the history,
updates the latest vehicle state, and sends invalid events to a separate place for review.

A second streaming task groups activity by zone and time window. It calculates active vehicle
counts, idle ratio, trips, earnings, and average speed. Counts are estimates because the live job
uses an approximate counting method and only sees the events received so far. Late events can also
fall outside the wait period.

An alert is made when a vehicle stays idle beyond the chosen limit. The system keeps the start of
the idle period so it does not create a new alert for the same idle period every time it checks.

\subsection{Daily batch processing}

Airflow looks for a new daily cost file and waits for the matching vehicle history. It starts a
Spark batch job once both inputs are available. The job reads all of the events for that simulated
day and the costs in the file.

The daily job calculates the final fare for each completed trip. It also calculates distance from
GPS points and estimates how much of the day each vehicle was working. It joins these figures with
the fuel and maintenance costs.

For each vehicle, the report shows income, costs, profit or loss, profit margin, distance, and
utilisation. The report can also show a missing cost record or a large difference between GPS
distance and the reported distance. Airflow checks that the result was created before marking the
day complete.

\subsection{Daily calculations}

\begin{table}[htbp]\centering\small
\caption{Main measures in the daily report.}
\begin{tabularx}{\textwidth}{@{}p{3.2cm}X@{}}
\toprule
\textbf{Measure} & \textbf{How it is found}\\
\midrule
Income & Add the final fare for each completed trip.\\
Distance & Estimate distance from consecutive GPS points. Ignore jumps that are too large to be real.\\
Utilisation & Find the share of vehicle updates that were not idle.\\
Total cost & Add the fuel and maintenance costs from the daily file.\\
Profit or loss & Subtract total cost from income.\\
Flags & Mark low profit, missing cost data, or a large difference in reported and GPS distance.\\
\bottomrule
\end{tabularx}
\end{table}

The daily result is the main report for a completed simulated day. The live view is useful during
the day, but it can be different because it is based on recent windows and partial data.

\section{Storage and Results for Users}

PostgreSQL holds the latest live vehicle information, recent zone summaries, alerts, and daily
profit results. The Parquet history keeps the full set of vehicle events. The daily report is also
saved as a CSV file and a web page.

The API reads the live and daily data and makes it available to the dashboard. The dashboard shows
which view is live and which view comes from the completed day. It also compares the two so users
can see how much the fast estimate differs from the daily result.

\imageplaceholder{dashboard.png}{The fleet dashboard showing current vehicle activity, daily results, and the comparison between the live and daily views.}{0.95}

\imageplaceholder{report.png}{The daily vehicle report showing income, costs, profit or loss, distance, and flags for issues that need review.}{0.95}

\section{Observability}

The project records structured logs, measures service health, and raises alerts. These features help
the team find out whether data is arriving, whether Spark is processing it, whether a daily run
finished, and whether the API can reach its database.

\subsection{Logs}

The services write structured log messages. Each message includes a time, service name, work stage,
event, and useful details. This makes it easier to follow an event from the source through
processing and storage.

\subsection{Measures and alerts}

Prometheus collects measures from the long-running services. The daily Spark job is short, so it
sends its measures to Pushgateway. Prometheus then reads those measures.

The measures cover incoming event rate, last event time, invalid event rate, Spark processing time,
processing delay, daily batch rows, API speed, and the age of the live and daily results. The
project includes fifteen alert rules. Examples include no recent vehicle events, a stopped Spark
job, too many invalid events, a late daily result, and a daily loss. A daily loss is a business
alert. It does not mean the monitoring system has failed.

Grafana shows the main measures in the Fleet Pipeline Health dashboard. A health check also confirms
that the API can reach PostgreSQL. Docker Compose uses service health checks to control start-up
order.

\imageplaceholder{grafana.png}{Grafana dashboard showing input rate, data freshness, invalid events, Spark work, daily batch results, and API health.}{0.95}

\section{Results}

During one verified run, the daily result for 17 August 2026 showed income of LKR 89,342.10 from
283 completed trips. The live view showed LKR 54,177.15 across 24 time windows. The live estimate
was LKR 35,164.95 lower, or 39.36 percent below the daily result.

\begin{table}[htbp]\centering\small
\caption{Example comparison for one simulated day.}
\begin{tabularx}{\textwidth}{@{}p{4.2cm}rX@{}}
\toprule
\textbf{View} & \textbf{Result} & \textbf{Meaning}\\
\midrule
Live earnings & LKR 54,177.15 & Estimated from incoming fare increases across 24 windows.\\
Daily income & LKR 89,342.10 & Calculated from 283 completed trips in the full day.\\
Difference & LKR -35,164.95 & The live estimate was lower than the daily result.\\
Difference as share of daily income & 39.36\% & The size of the gap in this run.\\
\bottomrule
\end{tabularx}
\end{table}

The largest known reason is the trip starting charge. The vehicle source adds LKR 120 to the
completed trip fare at the start of a trip. The live calculation adds fare increases during the
trip, but does not include this starting charge. For 283 trips, these charges total LKR 33,960.
After allowing for them, LKR 1,204.95 of the difference remains. This is about 1.35 percent of
the daily income. The rest is not assigned to one cause. The two paths use different time groups
and counting methods.

The trip counts also differ. A trip can appear in more than one live time or zone group. The live
count adds estimates from these groups. The daily job counts each completed trip once. This is why
the daily trip count is the better number for a full-day report.

These numbers are from one run. New results can change as the simulated clock moves forward.

\section{Limits and Future Improvements}

This is a learning project and has limits.

\begin{itemize}
  \item Spark runs in local mode. The project does not test how a large group of Spark machines
  handles a failure or shares work.
  \item Kafka uses one broker. A real service would use several brokers so data remains available
  if one machine stops.
  \item Events can be processed again after a restart. The database writes are made safe against
  repeats, but the Parquet history can contain duplicate rows.
  \item The Parquet data is not compacted. A long-running system would need to combine many small
  files or use a table format that manages this work.
  \item The expense source is simulated. Its costs are not a separate real bill from a fuel company
  or garage.
  \item The demo has no user login or encrypted connections. Those protections are needed before
  using real business or customer data.
\end{itemize}

At production scale, the team would use a managed or clustered Spark service, several Kafka
brokers, protected connections, access control, and alerts sent to an operator. The team would also
use a storage format that supports safe changes and cleanup of old data. The expense feed should
come from a real approved business source.

\section{Conclusion}

This project joins live vehicle data with a daily cost file. Kafka and Spark provide a current view.
Parquet keeps the history. Airflow starts the daily Spark job when both inputs are ready. PostgreSQL
and the API make the results available to the dashboard. Prometheus and Grafana help the team see
whether the system is healthy.

The Lambda design fits because the company needs two different answers. Dispatch needs a quick
view while vehicles are working. Finance needs a full-day result after costs arrive. The project
shows both answers and explains why they can be different.

\end{document}
